\chapter{仿真结果与参数依赖性分析}[Simulation Results and Parameter Dependence Analysis]

本章以“可对照实验的指标体系”为主线，给出端到端仿真的评估口径与数据组织方式。由于完整实验计划仍在迭代，本章强调“如何计算、如何验证、如何解释趋势”，具体数值结果将在后续批量实验中补齐。

\section{参数与实验设置}[Parameter Sets and Experimental Setup]

为便于后续对比与参数扫描，本节给出参考参数与可调范围。除原子谱线等物理常量外，其余参数均视为设计量，重点在于变化趋势而非单点数值。

\paragraph{原子系统参数}

\begin{table}[htbp]
\centering
\caption{$^{87}$Rb 原子系统参数}
\begin{tabular}{lcc}
\hline
参数 & 符号 & 数值 \\
\hline
跃迁波长 & $\lambda$ & 780 nm \\
自发辐射率 & $\gamma/2\pi$ & 6.07 MHz \\
自发辐射寿命 & $\tau = 1/\gamma$ & 26.2 ns \\
基态超精细分裂 & $\Delta_{\mathrm{hfs}}$ & 6.83 GHz \\
量子比特态 & $|0\rangle$, $|1\rangle$ & $|F=1, m_F=0\rangle$, $|F=2, m_F=0\rangle$ \\
\hline
\end{tabular}
\end{table}

该表用于标定发射的时间尺度与能级结构，是后续波包尺度与失谐设定的参考基准。

\paragraph{光学链路参数}

\begin{table}[htbp]
\centering
\caption{光学链路参考参数（可调）}
\begin{tabular}{lcc}
\hline
参数 & 符号 & 参考值 \\
\hline
QFC 转换效率 & $\eta_{\mathrm{QFC}}$ & 0.5 \\
QFC 噪声光子率 & $\mu_{\mathrm{noise}}$ & $10^{-3}$/bin \\
光纤损耗系数 & $\alpha$ & 0.2 dB/km \\
光纤长度（单臂） & $L$ & 25 km \\
光纤透过率 & $\eta_{\mathrm{fiber}}$ & 0.316 ($-5$ dB) \\
滤波器带宽 & $\Delta\nu_{\mathrm{filter}}$ & 1 GHz \\
\hline
\end{tabular}
\end{table}

这些参数用于刻画链路的主要损耗与噪声量级，参考相关实验报告与器件性能水平（如文献~\cite{vanleent2022entangling}），仿真中以此为基准进行趋势扫描。

\paragraph{探测器参数}

\begin{table}[htbp]
\centering
\caption{SNSPD 探测器参考参数（可调）}
\begin{tabular}{lcc}
\hline
参数 & 符号 & 参考值 \\
\hline
探测效率 & $\eta_d$ & 0.9 \\
暗计数率 & $R_{\mathrm{dark}}$ & 10 Hz \\
时间抖动 & $\sigma_t$ & 50 ps \\
死时间 & $\tau_{\mathrm{dead}}$ & 50 ns \\
\hline
\end{tabular}
\end{table}

探测器指标取自高性能 SNSPD 的典型范围，噪声部分在模型中统一并入有效噪声率，以便与实际计数率对应。

\paragraph{数值仿真参数}

\begin{table}[htbp]
\centering
\caption{数值仿真参数范围（可调）}
\begin{tabular}{lcc}
\hline
参数 & 符号 & 参考范围 \\
\hline
时间 bin 宽度 & $\Delta t$ & 0.5--1 ns \\
总 bin 数 & $N$ & 100--200 \\
Fock 空间截断 & $n_{\max}$ & 2--3 \\
MPS 键维度 & $\chi_{\max}$ & 50--64 \\
POVM 抽样次数 & $N_{\mathrm{shot}}$ & $10^2$--$10^5$ \\
\hline
\end{tabular}
\end{table}

该范围是精度与算力之间的折中：$\Delta t$ 决定时间分辨率，$N$ 决定波包覆盖范围，$n_{\max}$ 与 $\chi_{\max}$ 控制高光子数与纠缠截断误差。

\section{核心指标与验证}[Core Metrics and Verification]

端到端评估采用三类核心指标：两光子到达概率 $p_{\mathrm{arrive}}$、成功率的绝对与条件分解，以及成功样本的条件保真度 $F_{\mathrm{true}}$。具体而言，
\begin{align}
p_{\mathrm{success}}^{\mathrm{abs}} &= \sum_{m\in\mathcal{S}}\mathrm{Tr}(E_m\rho),\\
p_{\mathrm{success}}^{\mathrm{true}} &= \sum_{m\in\mathcal{S}}\mathrm{Tr}(E_m^{\mathrm{true}}\rho),\\
p_{\mathrm{success}}^{\mathrm{false}} &= p_{\mathrm{success}}^{\mathrm{abs}}-p_{\mathrm{success}}^{\mathrm{true}},\\
p_{\mathrm{success}}^{\mathrm{true}\mid\mathrm{arrive}} &= p_{\mathrm{success}}^{\mathrm{true}}/p_{\mathrm{arrive}}.
\end{align}
$p_{\mathrm{arrive}}$ 反映通道与器件的总损耗，$p_{\mathrm{success}}^{\mathrm{abs}}$ 给出每次尝试的实验可观测成功率，$p_{\mathrm{success}}^{\mathrm{false}}$ 刻画暗计数/背景导致的假成功占比，而 $F_{\mathrm{true}}$ 直接回答纠缠质量是否跨过经典极限。该口径与实验中“符合门控后统计”一致，也便于把“到达损失”和“测量误判”分开讨论。

\paragraph{单臂发射验证}
发射波包由驱动脉冲 $\Omega(t)$ 与腔参数给出的有效耦合率 $\gamma_{\mathrm{eff}}$ 共同决定，离散门通过 $H_{\mathrm{sys}}(t_n)$ 与 $L_\mu$ 直接生成波包。这里主要检查三件事：波包形状是否随驱动改变而合理移动/展宽，按 bin 统计的非真空概率之和是否与预期发射概率一致，以及 $\langle \hat{n}(t_n) \rangle$ 的形状是否与驱动包络吻合。图 \ref{fig:wavepacket} 给出示例。为控制数值成本与精度，进一步观察 MPS 键维度随时间变化（图 \ref{fig:bond_dim}），并据此选择 $\chi_{\max}$。

\begin{figure}[htbp]
\centering
\includegraphics[width=0.85\linewidth]{fig_wavepacket}
\caption{发射波包时间结构}
\label{fig:wavepacket}
\end{figure}

\begin{figure}[htbp]
\centering
\includegraphics[width=0.85\linewidth]{fig_bond_dim}
\caption{发射过程中 MPS 键维度的演化}
\label{fig:bond_dim}
\end{figure}

\paragraph{单光子到达时间直方图与接受窗优化}
实验上最直接的标定量之一是“到达时间直方图 + 接受窗口”。在仿真中以点击记录的 bin 分布构造单臂时间直方图，并在其上叠加接受窗口宽度 $\Delta t_{\mathrm{win}}$。对每个窗口宽度扫描输出：
单臂/双臂的信噪比（SBR）、窗口接受率（接受事件占比）、$p_{\mathrm{success}}^{\mathrm{abs}}$ 与 $F_{\mathrm{true}}$ 的变化趋势。该扫描给出“速率—保真度”在窗口选择上的直接权衡，是对实验调参最具指导意义的图之一。

\begin{figure}[htbp]
\centering
\includegraphics[width=0.85\linewidth]{fig_accept_hist}
\caption{单臂点击时间直方图与接受窗口示意（含背景噪声）。}
\label{fig:accept_hist}
\end{figure}

\begin{figure}[htbp]
\centering
\includegraphics[width=0.85\linewidth]{fig_window_tradeoff}
\caption{接受窗口宽度扫描：接受率、SBR、$F_{\mathrm{true}}$ 与假成功占比的权衡示意。}
\label{fig:window_tradeoff}
\end{figure}

\paragraph{双臂纠缠与 HOM 验证}
理想情况下，BSM 成功概率为
\begin{equation}
P_{\mathrm{success}}^{\mathrm{ideal}} = \frac{1}{2}|\langle\phi_A|\phi_B\rangle|^2,
\end{equation}
其中 $1/2$ 来自线性光学 BSM 的本征效率限制。实际链路中需同时报告 $p_{\mathrm{arrive}}$ 与 $F_{\mathrm{true}}$，以区分损耗与误判。HOM 干涉通过扫描两臂到达时延得到
\begin{equation}
P_{\mathrm{coinc}}(\delta t) = \frac{1}{2}\left[1 - V \cdot e^{-|\delta t|/\tau_c}\right],
\end{equation}
图 \ref{fig:hom} 给出示例曲线（通常以 wrong coincidences 归一化计数表示）。若 dip 的宽度或中心位置与单臂波包统计不一致，通常意味着时延标定、相位漂移或偏振失配未被充分补偿；可见度下降则主要来自背景噪声、时间抖动与模式不匹配。

\begin{figure}[htbp]
\centering
\includegraphics[width=0.85\linewidth]{fig_hom}
\caption{HOM 干涉曲线}
\label{fig:hom}
\end{figure}

\paragraph{BSM 点击模式与真假事件拆分}
中心站四探测器的点击模式统计是检验模型可信度的关键。仿真中输出每个双击模式的概率分布（成功模式与失败模式分开），并用暗计数拆分给出
\begin{equation}
f_{\mathrm{false}}=\frac{p_{\mathrm{success}}^{\mathrm{false}}}{p_{\mathrm{success}}^{\mathrm{abs}}},
\end{equation}
从而直接评估“成功事件里有多少是暗计数/背景伪造”。该统计可随链路长度、接受窗口与背景噪声扫描，形成对实验最直观的误判预算。

\begin{figure}[htbp]
\centering
\includegraphics[width=0.85\linewidth]{fig_bsm_patterns}
\caption{中心站双击模式分布与暗计数拆分示意。}
\label{fig:bsm_patterns}
\end{figure}

\paragraph{原子关联曲线与 CHSH 指标}
对纠缠质量的实验表述通常采用三基相关曲线（X/Y/Z）或 CHSH 不等式。仿真中可对后验原子态计算相关函数 $E(\theta_A,\theta_B)$ 并给出
\begin{equation}
S = E(\theta_A,\theta_B)-E(\theta_A,\theta_B')+E(\theta_A',\theta_B)+E(\theta_A',\theta_B'),
\end{equation}
以量化纠缠是否超出经典界 $S>2$。在存在泄漏态或读出误差时，可同时给出“物理态 fidelity”与“可观测 fidelity”的分离结果，为实验读出提供对照。

\begin{figure}[htbp]
\centering
\includegraphics[width=0.85\linewidth]{fig_chsh}
\caption{CHSH 量 $S$ 随链路长度变化的示意曲线（虚线为经典上界 $S=2$）。}
\label{fig:chsh}
\end{figure}

\section{参数依赖与误差来源}[Parameter Dependence and Error Sources]

参数扫描以“趋势判断”为主：改变某一物理量并观察 $p_{\mathrm{arrive}}$、$P(\mathrm{success}\,|\,\mathrm{arrive})$ 与 $F_{\mathrm{true}}$ 的系统性变化，同时用误差预算框架定位主导来源。

\paragraph{QFC 噪声}
保真度随噪声率 $\mu_{\mathrm{noise}}$ 单调下降，低噪声区间近似线性退化：
\begin{equation}
F(\mu_{\mathrm{noise}}) \approx F_0 - \alpha_{\mathrm{noise}}\mu_{\mathrm{noise}}.
\end{equation}
图 \ref{fig:qfc_noise} 给出趋势示例。该趋势反映了“假成功事件”占比上升对条件态的直接污染。
在这一过程中，$p_{\mathrm{arrive}}$ 主要由链路损耗决定，变化不大；下降更明显的是 $p_{\mathrm{success}}^{\mathrm{true}\mid\mathrm{arrive}}$ 与 $F_{\mathrm{true}}$，而 $f_{\mathrm{false}}$（假成功占比）随噪声显著上升，直接反映噪声对宣告成功的稀释效应。

\begin{figure}[htbp]
\centering
\includegraphics[width=0.85\linewidth]{fig_qfc_noise}
\caption{保真度随噪声率变化}
\label{fig:qfc_noise}
\end{figure}

\paragraph{光纤长度}
成功率随距离指数衰减：
\begin{equation}
p_{\mathrm{success}}^{\mathrm{abs}}(L) = p_{\mathrm{success}}^{\mathrm{abs}}(0)\times 10^{-\alpha L/5}.
\end{equation}
图 \ref{fig:distance} 展示成功率与保真度随距离的变化。远距离下噪声占比上升，保真度出现拐点，这是链路极限的主要判据之一。
从指标分解上看，$p_{\mathrm{arrive}}$ 会随距离快速下降，而 $F_{\mathrm{true}}$ 在中短距离内相对稳定，直到噪声与信号可比时才明显下降；此时 $f_{\mathrm{false}}$ 的上升成为性能转折的主要信号。

\begin{figure}[htbp]
\centering
\includegraphics[width=0.85\linewidth]{fig_distance}
\caption{成功率与保真度随光纤长度变化}
\label{fig:distance}
\end{figure}

\paragraph{探测器效率与时间分辨}
探测效率提升使成功率近似按 $\eta_d^2$ 增长，同时改善信噪比；时间 bin 宽度过大将削弱时间分辨与 HOM 可见度，过小则提高噪声占比与计算成本。图 \ref{fig:detector_eff} 给出探测效率的趋势示例。
实际设置中，需要同时考虑“时间窗覆盖波包”与“噪声积分窗口”的平衡：窗口越窄越利于保真度，但会牺牲成功率。

\begin{figure}[htbp]
\centering
\includegraphics[width=0.85\linewidth]{fig_detector_eff}
\caption{探测器效率对成功率与保真度的影响}
\label{fig:detector_eff}
\end{figure}

\paragraph{误差预算与损耗分解}
保真度损失可写为
\begin{equation}
1 - F = \epsilon_{\mathrm{QFC}} + \epsilon_{\mathrm{dark}} + \epsilon_{\mathrm{PMD}} + \epsilon_{\mathrm{mismatch}} + \epsilon_{\mathrm{other}},
\end{equation}
成功率损失可分解为
\begin{equation}
\frac{p_{\mathrm{success}}^{\mathrm{abs}}}{p_{\mathrm{success}}^{\mathrm{ideal}}} = \eta_{\mathrm{QFC}}^2 \times \eta_{\mathrm{fiber}}^2 \times \eta_d^2 \times \eta_{\mathrm{other}}.
\end{equation}
通过逐一关闭或单独开启噪声源并比较 $\Delta F$ 与 $f_{\mathrm{false}}$，即可建立误差预算并识别主导瓶颈。

\section{优化策略与取舍}[Optimization and Trade-offs]

提升保真度的关键路径是降低背景噪声（QFC 噪声与暗计数）、缩小符合窗口并提升偏振/相位稳定性；提升成功率则依赖高效率 QFC、低损耗光纤与复用策略。二者间的取舍主要由窗口选择与实验目标决定：宽窗口提高成功率但引入噪声，窄窗口提升保真度但牺牲速率。实际应用中可用“rate–fidelity 前沿曲线”来选工作点：横轴为事件率（$p_{\mathrm{success}}^{\mathrm{abs}}\times$重复频率），纵轴为 $F_{\mathrm{true}}$ 或 CHSH $S$。

\begin{figure}[htbp]
\centering
\includegraphics[width=0.85\linewidth]{fig_rate_fidelity}
\caption{事件率–保真度前沿示意（不同工作点的权衡曲线）。}
\label{fig:rate_fidelity}
\end{figure}

更实际的做法是先判断系统处于“损耗主导”还是“噪声主导”：若 $p_{\mathrm{arrive}}$ 极低且 $f_{\mathrm{false}}$ 仍小，则应优先提高收集与传输效率；若 $f_{\mathrm{false}}$ 已接近或超过信号，则应优先压噪并收紧判据。该判断可直接从 $p_{\mathrm{arrive}}$、$p_{\mathrm{success}}^{\mathrm{true}\mid\mathrm{arrive}}$ 与 $F_{\mathrm{true}}$ 的相对变化中读出。

\section{本章小结}[Chapter Summary]

本章给出端到端仿真的指标体系与实验设计思路，强调以有效事件归一化的统计口径来分离“到达损失”和“测量判据”的影响。通过参数趋势与误差预算，可在尚未固定最终实验参数时识别关键瓶颈与优化方向。后续将结合批量实验数据补充定量曲线与对照表格，使分析与实验需求直接对齐。
